\documentclass[10pt,a4paper]{amsart}
\usepackage[margin=19mm]{geometry}
\usepackage{amsmath,amssymb,mathtools}
\usepackage[scr=rsfs,cal=euler]{mathalfa}
\usepackage{xfrac}
\usepackage[unicode,hidelinks,bookmarks=false]{hyperref}
\newcommand{\ie}{i.e.\ }
\newcommand{\cf}{cf.\ }
\newcommand{\resp}{respectively\ }
\newcommand{\Subsection}{\subsection}
\providecommand{\nobreakdash}{\nobreak\hbox{-}}
\setlength{\emergencystretch}{2em}
\multlinegap0pt
\usepackage{bm}
\usepackage{bbm}
\usepackage{dsfont}
\usepackage{xspace}
\usepackage{cancel}
\usepackage{mathtools}
\usepackage{amsthm}
\makeatletter
\@ifundefined{theorem}{\newtheorem{theorem}{Theorem}}{}
\@ifundefined{lemma}{\newtheorem{lemma}[theorem]{Lemma}}{}
\@ifundefined{remark}{\newtheorem{remark}[theorem]{Remark}}{}
\makeatother
\title[Orthonormal Strichartz estimates for Dunkl-Schr\"{o}dinger e]{Orthonormal Strichartz estimates for Dunkl-Schr\"{o}dinger equation of initial data with Sobolev regularity}
\author{Guoxia Feng, Shyam Swarup Mondal, Manli Song, Huoxiong Wu}
\date{Source: arXiv:2506.05493v2 (CC BY 4.0)}
\begin{document}
\maketitle

\noindent\emph{Source and licence.} Orthonormal Strichartz estimates for Dunkl-Schr\"{o}dinger equation of initial data with Sobolev regularity. Version arXiv:2506.05493v2 (CC BY 4.0), distributed under the Creative Commons Attribution 4.0 International licence, as recorded in the arXiv licence metadata for this exact version and shown on the abstract page https://arxiv.org/abs/2506.05493v2. Full rights notice: licenses/arxiv-2506.05493-CC-BY-4.0.txt.\par\smallskip

\noindent\textit{Editorial scope.} Counted unit: the Theorem whose source label is \texttt{adimissible} in the article (source0.tex lines 921--927, source label \texttt{adimissible}), delivered together with the article's own complete original proof of it (source0.tex lines 928--957, one \texttt{proof} environment of 30 lines). No mathematics is written, repaired or re-derived: statement and proof are the article's own text line for line. Required context retained uncounted: translator (lines 541--543); Young (lines 557--559); Schrodinger (lines 568--570); cut-off decay (lines 593--599); duality-principle (lines 668--677); Lorentz-refined (lines 867--873). Every definition, display and label of those lines is retained unchanged. Omitted from this package and not redistributed: 1--540, 544--556, 560--567, 571--592, 600--667, 678--866, 874--920, 958--1418. Nothing inside a retained excerpt is omitted or altered: every retained body is delivered exactly as the article's lines are written, so no omission receipt of any kind accompanies this package. Citations: the retained excerpts carry no citation command at all, so no bibliography entry of the article is transcribed and no reference is redistributed. Reference adaptation: none was needed and none was made; every reference inside the delivered unit resolves inside it, so no unresolved reference or citation is produced. Printed slips kept literal: no printed word, symbol, hypothesis or number of the counted statement or of its proof was altered. Layout and class: the article's own class is \texttt{amsart} with packages including amssymb, amsthm, amsmath, enumerate, ifthen, graphicx, cite, tikz, amscd, latexsym, hyperref; the wrapper is the lane's pinned \texttt{amsart} editorial wrapper at 10pt on A4, which restates the theorem-like environments the retained text needs and adds the article's own macro definitions that the retained text needs (none), with extra wrapper packages (bm, bbm, dsfont, xspace, cancel, mathtools). The delivered numbering is the unit's own. Assets: no retained span contains a figure environment, a graphics-inclusion command, a PDF-page inclusion, a table or a TikZ picture; no image file is copied into this package, the package declares no asset dependency and no image file is redistributed with it. Licence: version arXiv:2506.05493v2 of this article is distributed under the Creative Commons Attribution 4.0 International licence, as recorded in the arXiv licence metadata for this exact version and shown on the abstract page https://arxiv.org/abs/2506.05493v2. A full rights notice is delivered at licenses/arxiv-2506.05493-CC-BY-4.0.txt. Characters: every retained excerpt is pure ASCII, so the delivered engine, XeLaTeX with its default Latin Modern text font, reports no missing character.
\begin{equation} \label{translator}
\tau_xf(y)=V_\kappa[\tilde{f}\left(\sqrt{|y|^2+|x|^2-2\langle y,\cdot\rangle}\right)](x)=\int_{\mathbb{R}^n}\tilde{f}\left(\sqrt{|y|^2+|x|^2-2\langle y,\xi\rangle}\right)d\rho^\kappa_x(\xi).
\end{equation}

\begin{equation}\label{Young}
\|f*_\kappa g\|_{L_\kappa^r(\mathbb{R}^n)}\leq \|f\|_{L_\kappa^p(\mathbb{R}^n)} \|g\|_{L_\kappa^q(\mathbb{R}^n)}.
\end{equation}

\begin{equation}\label{Schrodinger}
K_{it}(x)=\mathcal{F}^{-1}_\kappa\left(e^{-it|\cdot|^2}\right)(x)=\frac{1}{c_\kappa |t|^\frac{N}{2}}e^{-i\frac{N\pi}{4} sgn t}e^{i\frac{|x|^2}{4t}}.
\end{equation}

\begin{remark}
The Dunkl kernel of the localised  Dunkl-Schr\"{o}dinger operator $e^{it\Delta_\kappa}P$ is also radial and given by $\mathcal{F}^{-1}_\kappa (e^{-it|\cdot|^2} \psi(|\cdot|))$ satisfying the following dispersive estimate, i.e.,
\begin{equation}\label{cut-off decay}
   \left|\tau_x\mathcal{F}^{-1}_\kappa \left(e^{-it|\cdot|^2} \psi(|\cdot|)\right)(y)\right|\lesssim_\psi |t|^{-\frac{N}{2}}, \quad \forall x,y\in\mathbb{R}^n,
\end{equation}
which follows from Young's inequality \eqref{Young}, \eqref{translator} and \eqref{Schrodinger}.
\end{remark}

 \begin{lemma}[Duality principle]\label{duality-principle} Assume that $A$ is a bounded linear operator from $L_\kappa^2(\mathbb{R}^n)$ to $L^{2q,2r}(\mathbb{R}, L_\kappa^{2p}(\mathbb{R}^n))$ for some $p,q,\alpha\geq 1$. Then 
         \begin{equation*}
			\bigg\|\sum_{j=1}^\infty \lambda_j|A f_j|^2\bigg\|_{L^{q,r}(\mathbb{R}, L_\kappa^{p}(\mathbb{R}^n))}\lesssim \|\{\lambda_j\}^\infty_{j=1}\|_{\ell^{\alpha}} ,
		\end{equation*}	
  holds for all families of orthonormal functions $\{f_j\}_{j=1}^\infty$ in $L_\kappa^2(\mathbb{R}^n)$ and all sequences $\{\lambda_j\}^\infty_{j=1}$ in $\ell^{\alpha}$ if and only if
		\begin{equation*}
	\big\|WAA^*\overline{W}\big\|_{\mathfrak{S}^{\alpha^\prime}(L^2(\mathbb{R}, L_\kappa^2(\mathbb{R}^n)))}\lesssim \big\|W\big\|^2_{L^{2q^\prime,2r^\prime}(\mathbb{R}, L_\kappa^{2p^\prime}(\mathbb{R}^n))},
		\end{equation*}
  for all $W\in L^{2q^\prime,2r^\prime}(\mathbb{R}, L_\kappa^{2p^\prime}(\mathbb{R}^n))$, where the function $W$ is interpreted as an operator which acts by multiplication.
	\end{lemma}

\begin{theorem} \label{Lorentz-refined}
If $N\geq 1$, $(\frac{1}{p},\frac{1}{q})\in (A,F]$ and $\alpha=\frac{2p}{p+1}$, then we have
\begin{equation*}
			\bigg\|\sum_{j=1}^\infty \lambda_j|e^{it\Delta_\kappa}f_j|^2\bigg\|_{L^{q,\alpha}(\mathbb{R}, L_\kappa^p(\mathbb{R}^n))}\lesssim \|\{\lambda_j\}^\infty_{j=1}\|_{\ell^{\alpha}} 
\end{equation*}	
  holds for all families of orthonormal functions $\{f_j\}_{j=1}^\infty$ in $L_\kappa^2(\mathbb{R}^n)$ and all sequences $\{\lambda_j\}^\infty_{j=1}$ in $\ell^{\alpha}$. 
\end{theorem}

\begin{theorem} \label{adimissible}
If $N\geq 1$, $(\frac{1}{p},\frac{1}{q})\in [B,A)$ and $\alpha=\frac{2p}{p+1}$, then we have
\begin{equation*}
			\bigg\|\sum_{j=1}^\infty \lambda_j|e^{it\Delta_\kappa}Pf_j|^2\bigg\|_{L^q(\mathbb{R}, L_\kappa^p(\mathbb{R}^n))}\lesssim_\psi \|\{\lambda_j\}^\infty_{j=1}\|_{\ell^{\alpha}} 
\end{equation*}	
  holds for all families of orthonormal functions $\{f_j\}_{j=1}^\infty$ in $L_\kappa^2(\mathbb{R}^n)$ and all sequences $\{\lambda_j\}^\infty_{j=1}$ in $\ell^{\alpha}$. 
\end{theorem}

\begin{proof}
    From the fact that $e^{it\Delta_\kappa}$ is unitary on $L_\kappa^2(\mathbb{R}^n)$ and $\|Pf_j\|_{L_\kappa^2(\mathbb{R}^n)}\leq \|f_j\|_{L_\kappa^2(\mathbb{R}^n)}$, it immediately follows that
    \begin{equation*}
			\bigg\|\sum_{j=1}^\infty \lambda_j|e^{it\Delta_\kappa}Pf_j|^2\bigg\|_{L^\infty(\mathbb{R}, L_\kappa^1(\mathbb{R}^n))}\leq \sum_{j=1}^\infty |\lambda_j|\|e^{it\Delta_\kappa}Pf_j\|^2_{L^\infty(\mathbb{R}, L_\kappa^2(\mathbb{R}^n))}\leq \|\{\lambda_j\}^\infty_{j=1}\|_{\ell^1}. 
\end{equation*}	
By interpolation, it suffices to prove that 
\begin{equation*}
			\bigg\|\sum_{j=1}^\infty \lambda_j|e^{it\Delta_\kappa}Pf_j|^2\bigg\|_{L^q(\mathbb{R}, L_\kappa^p(\mathbb{R}^n))}\lesssim \|\{\lambda_j\}^\infty_{j=1}\|_{\ell^{\alpha}} 
\end{equation*}	
holds for $1+\frac{2}{N}\leq p<\frac{N+1}{N-1}$. By duality principle  given in Theorem \ref{duality-principle}, it is equivalent to  prove
\begin{equation}\label{P-T_1}
	\big\|We^{it\Delta_\kappa}(e^{it\Delta_\kappa})^*P^2\overline{W}\big\|_{\mathfrak{S}^{\alpha^\prime}(L^2(\mathbb{R}, L_\kappa^2(\mathbb{R}^n)))}\lesssim_\psi\big\|W\big\|^2_{L^{2q^\prime}(\mathbb{R}, L_\kappa^{\alpha^\prime}(\mathbb{R}^n))}
		\end{equation}
  for all $W\in L^{2q^\prime}(\mathbb{R}, L_\kappa^{\alpha^\prime}(\mathbb{R}^n))$, where the operator $e^{it\Delta_\kappa}(e^{it\Delta_\kappa})^*P^2$ is given by
  \begin{equation*}
      e^{it\Delta_\kappa}(e^{it\Delta_\kappa})^*P^2g(x)=\int_\mathbb{R} e^{i(t-s)\Delta_\kappa} P^2g(\cdot,s)(x)ds.
  \end{equation*}
 Similarly to Theorem \ref{Lorentz-refined},  consider the analytic family of operators $\mathcal{T}_z$ defined by
\begin{equation*}
    \mathcal{F}_{t,\kappa}(\mathcal{T}_zg)(\xi,\tau)=\frac{1}{\Gamma(z+1)}(\tau-|\xi|^2)_+^z\psi^2(|\xi|)\mathcal{F}_{t,\kappa}g(\xi,\tau)
\end{equation*}
for $(\xi,\tau)\in\mathbb{R}^n\times\mathbb{R}$. Then the kernel of $\mathcal{T}_z$ is given by
\begin{align*}
    &\quad \tau_x\left(\frac{1}{c_\kappa\Gamma(z+1)}\int_{\mathbb{R}}\tau_+^ze^{i\tau (t-s)}d\tau\int_{\mathbb{R}^n} e^{i(t-s)|\xi|^2} \psi^2(|\xi|)K(-y,i \xi)h_\kappa(\xi)d\xi\right)\\
    &=ie^{iz \pi/2}(t-s+i0)^{-z-1}\tau_x\mathcal{F}^{-1}_\kappa \left(e^{i(t-s)|\cdot|^2} \psi^2(|\cdot|)\right)(-y)
\end{align*}
and by \eqref{cut-off decay}, the kernel of $W_1\mathcal{T}_z W_2$ is also bounded by
$$C_\psi(Im(z))|W_1(x,t)||t-s|^{-Re(z)-1-\frac{N}{2}}|W_2(y,s)|.$$
Then, proceeding similarly as in  Theorem \ref{Lorentz-refined}, we have our required estimate \eqref{P-T_1}. 
\end{proof}

\end{document}
